This section defines the architectural role of repeated scalar execution in Bedrock. The architecture exposes
scalar register state and bounded repeated instruction forms. A high-performance implementation may recognize
regular repeated operations and execute them with wider internal datapaths, but that internal width is not
architectural state.

\subsection{Architectural Scalar Semantics}
REPcc and REPG are architecturally scalar. Their visible behavior is the same as executing the repeated
instruction or byte-counted instruction group in program order under the selected counter and termination rules.
The architectural counter, FLAGS image, memory effects, and auto-update register effects are observed as if each
committed iteration completed before the next one began.

\begin{itemize}
\item register and memory operands are evaluated according to the repeated instruction's normal rules;
\item auto-update effective-address forms update as if each iteration executed the scalar instruction;
\item faults and exceptions are precise at the repeated-operation boundary defined by the REPcc or REPG rule;
\item FLAGS and FFLAGS are updated according to the repeated instruction and repeat accumulation rules;
\item software must not depend on an implementation's internal vector width or streaming datapath width.
\end{itemize}

\subsection{Implementation Freedom}
An implementation may recognize regular repeated operations and execute them with a micro-loop, loop buffer,
streaming engine, or wider internal datapath. That internal width is not architectural state and is not part of
the ABI. This freedom does not permit reordering visible exceptions, changing the committed counter value, or
exposing partial updates within an iteration.

An implementation may also choose scalar execution for any repeated operation. A streaming candidate is a shape
that can be recognized; it is not a throughput, latency, or vector-width guarantee.

\subsection{Streaming Candidates}
Regular memory traversal forms such as MOV, CMP, EXTZ, EXTS, cache maintenance, and selected atomic-free
transforms are natural streaming candidates when combined with REPcc or REPG. An implementation may also stream
fixed-size instruction groups when REPG describes a bounded group that has no unresolved control transfer into or
out of the group body.

@STREAMING_CANDIDATE_TABLE@

\subsection{Compiler Model}
Compilers may emit repeated scalar forms for compact code size and predictable restart state. They do not need to
know whether a specific processor executes the loop with scalar issue or a streaming engine. Generated code must
still preserve aliasing, ordering, and exception behavior exactly as required by the scalar instruction sequence.

\subsection{Relationship to Other Acceleration}
REPG uses an unsigned byte count for the grouped body. The byte count must describe a nonempty body that ends on
an instruction boundary. The counter update occurs only after a whole group iteration succeeds.

Faults remain precise at the repeated-operation boundary. Completed iterations are committed; an incomplete
faulting iteration is reported at the faulting instruction with restart state sufficient to continue or emulate
the remaining work.

@REPEAT_RESTART_TABLE@
